\subsection{Field Notes: The 4-Layer Epistemic Standard}

Raw notes must strictly separate what was observed from what was inferred. Every interview record must maintain 4 distinct layers:
\begin{enumerate}[leftmargin=1.2em,itemsep=0.2ex]
  \item \textbf{Layer 1: Verbatim Speech.} Exact words and technical vocabulary used by the participant, enclosed in quotation marks.
  \item \textbf{Layer 2: Observable Actions and Artifacts.} Measurable facts: lines of code audited, tools open on screen, dollars spent on invoices, dates of release delays.
  \item \textbf{Layer 3: Founder Interpretation.} The researcher's provisional hypotheses regarding underlying jobs, anxieties, and unstated motives.
  \item \textbf{Layer 4: Counter-Hypotheses and Unresolved Questions.} Facts that contradict the founder's thesis and questions requiring follow-up.
\end{enumerate}

Synthesize notes within 6 hours of the conversation. Delayed synthesis results in human memory flattening uncomfortable contradictions and exaggerating favorable comments.

\subsection{Cross-Case Synthesis vs. Mention Counting}

Novice teams tally keywords: ``8 respondents asked for mentorship, 5 asked for certifications.'' Word counting is intellectually lazy: 8 respondents use the word ``mentorship'' to mean 8 completely different things (from casual career cheerleading to line-by-line algorithmic debugging).

Deploy \textbf{Cross-Case Mechanism Synthesis}:
\begin{itemize}[leftmargin=1.2em,itemsep=0.2ex]
  \item Construct an empirical matrix with individual respondents as rows and structural decision variables as columns: Triggering Incident, Baseline Capability, Existing Workaround, Budget Owner, Critical Veto, and Dominant Anxiety.
  \item Group rows by the \textbf{underlying causal mechanism} that explains their choice.
  \item Actively highlight contradictions: Why did Respondent 4 switch to an external consultant while Respondent 7, in an identical role, solved the issue with open-source documentation?
\end{itemize}

\subsection{Theoretical Saturation Is Local and Conditional}

Founders often proclaim: ``We hit saturation after 12 interviews because we stopped hearing new things.'' True saturation is always conditional on the sampling boundary. 12 interviews with backend engineers at venture-funded startups in Bangalore tell you nothing about engineering teams at enterprise banks in Mumbai or academic departments in Chennai.

A disciplined stopping rule for an exploratory cycle:
\begin{itemize}[leftmargin=1.2em,itemsep=0.2ex]
  \item Within this tightly defined segment, can we reliably predict the next interviewee's workaround, dominant alternative, and primary switching barrier?
  \item Have we actively sought out and interviewed negative, disconfirming cases?
  \item Do we possess sufficient behavioral clarity to design a bounded, paid commercial test?
\end{itemize}
When these 3 conditions are met, cease exploratory interviewing and execute a behavioral transaction test.

\subsection{Surveys Answer Different Questions Than Interviews}

Interviews excel at discovering unknown mechanisms, unexpected workarounds, and nuanced emotional trade-offs. They cannot measure population prevalence or statistical distribution. Quantitative surveys, when properly designed, measure frequency, proportions, and demographic correlations across a defined population [R12].

However, a poorly designed survey sent to a self-selected email newsletter or Twitter following is worse than useless: it yields statistically invalid numbers coated in a veneer of pseudo-scientific authority. Stated preferences on surveys suffer from massive hypothetical bias: respondents rate features as ``extremely important'' when they bear zero cost for checking the box [R13].

\subsection{Designing Surveys Around Past Decisions, Not Applause}

If a quantitative survey is deployed, anchor questions strictly on verifiable past behavior:
\begin{itemize}[leftmargin=1.2em,itemsep=0.2ex]
  \item Bad: ``On a scale of 1 to 5, how interested are you in mastering production AI systems?''
  \item Good: ``In the past 12 months, which of the following actions have you personally taken to learn AI systems engineering? (Select all that apply: Purchased a paid book; Completed a paid course; Attended a corporate workshop; Built an open-source project; None of the above).''
  \item Bad: ``Would you pay ₹50,000 for a course that guaranteed expert feedback?''
  \item Good: ``What was the total financial expenditure you or your employer approved for technical upskilling in the last fiscal year? (₹0; Under ₹10,000; ₹10,000 to ₹50,000; Above ₹50,000).''
\end{itemize}
Pretest surveys with cognitive interviews to ensure respondents interpret questions as intended before launching to large samples.
